\documentclass[11pt]{article}
\usepackage[margin=1in]{geometry}
\usepackage{amsmath,amssymb,amsthm}
\usepackage{booktabs}
\usepackage[hidelinks]{hyperref}
\usepackage{url}

\newtheorem{lemma}{Lemma}
\newtheorem{theorem}{Theorem}

\title{Interval-certified bounds on the area of the Mandelbrot set}
\author{Bj\"orn Lindheim}
\date{Version 2, 11 November 2026 \\[0.5ex] \small Revised after an independent review; see the revision note at the end.}

\begin{document}
\maketitle

\begin{abstract}
Using ball-arithmetic certificates checked by separately implemented verifiers, we establish that the planar Lebesgue area of the Mandelbrot set satisfies
\[
  1.5064820 \;\le\; \mu(M) \;\le\; 6880083171681/2^{42} \;<\; 1.5643498 .
\]
The upper bound comes from a depth-17 dyadic quadtree (with two further levels of exterior refinement) whose exterior cells are certified by escape or by a Koebe-type distance estimate with explicit truncation bounds; it is below Fisher and Hill's 1993 bound and below the certified upper bounds we found, but above F\"orstemann's uncertified 1.53121. The lower bound sums certified lower bounds on the areas of 3,352,592 hyperbolic components of periods 3 to 504, plus the exact areas of the two components of periods 1 and 2. A stronger ball-arithmetic lower bound, $\mu(M)>1.50651$, was published in October 2026 by the same method; ours is an independent certification by that method and is weaker. All arithmetic is Arb ball arithmetic via python-flint; nothing is formalised in a proof assistant.
\end{abstract}

\section{Introduction}

The area of the Mandelbrot set $M=\{c\in\mathbb{C}: \text{the orbit of }0\text{ under }z\mapsto z^2+c\text{ is bounded}\}$ is not known in closed form. Pixel counting gives $\mu(M)\approx1.5065918849\pm2.8\cdot10^{-9}$ (F\"orstemann 2012, via \cite{munafo}). Fisher and Hill \cite{fisherhill1993} gave $1.50296686<\mu(M)<1.57012937$ from a quadtree and the Koebe quarter theorem, computed in floating point (we quote these values and the date as they are usually cited; we have not resolved the bibliographic details of the original manuscript). Hill then searched for centres of hyperbolic components, followed recursively attached components and evaluated component areas, obtaining lower bounds of $1.50585063$ (1996) and $1.506303622$ (1997, 430,809 components) \cite{munafo,hill2003}. F\"orstemann \cite{forstemann2017} reports $1.50640<\mu(M)<1.53121$ from a depth-26 quadtree in double precision. These computations do not control floating-point error. Upper bounds from the Laurent series of the B\"ottcher map \cite{ewingschober1992,bittner2017,irving_mandelbrot} converge slowly; the best we found is $\mu(M)\le1.651587$ \cite{irving_mandelbrot}.

Interval-certified results: Heiland-Allen reports proven bounds $1.416492462158203125\le\mu(M)\le1.84781646728515625$ \cite{heilandallen2023}. Irving's verified rendering, part of a Lean formalisation of the Mandelbrot set \cite{irving_ray,irving_render}, is quoted as giving $\mu(M)\le1.826454$ at depth 12; we have not independently checked that figure. A study published on 5 October 2026 \cite{sciencejournal2026} proves $\mu(M)>1.50651$ (reported as $1.5065108996$) in ball arithmetic by summing area-theorem lower bounds over 1,234,827 hyperbolic components; it gives no upper bound.

\begin{theorem}\label{thm:main}
$1.5064820\le\mu(M)\le 6880083171681/2^{42}$, where $6880083171681/2^{42}=1.5643497981002383\ldots$
\end{theorem}

The methods are Hill's (component sums) and Fisher and Hill's (quadtree with a distance estimate). The contribution of this note is the explicit two-sided interval with its certificates and their independent verification: the upper bound is, as far as we have found, the best certified upper bound, while the lower bound is weaker than \cite{sciencejournal2026}. We make no claim of priority.

\section{The lower bound: hyperbolic components}

A hyperbolic component $W$ of period $p$ is a connected component of the set of $c$ for which $z\mapsto z^2+c$ has an attracting cycle of exact period $p$. Each $W$ has a unique centre $c^*$, where $0$ lies on the cycle, and the multiplier map $\lambda:W\to\mathbb{D}$ is a conformal isomorphism with $\lambda(c^*)=0$; this holds for primitive and satellite components alike \cite[Thm.~6.5]{milnor2000}, see also \cite{douadyhubbard}. Distinct components are disjoint open subsets of $\operatorname{int}M$.

\begin{lemma}[component area]\label{lem:area}
Let $\psi(\lambda)=c^*+\sum_{n\ge1}a_n\lambda^n$ be the inverse of the multiplier map of $W$. Then $\mu(W)=\pi\sum_{n\ge1}n|a_n|^2$, so any finite partial sum is a lower bound.
\end{lemma}
\begin{proof}
For $0<r<1$, $\psi$ is injective on $D_r$, so $\mu(\psi(D_r))=\int_{D_r}|\psi'|^2\,dA=\pi\sum_n n|a_n|^2r^{2n}$ by orthogonality of $\lambda^n$ on circles. Let $r\uparrow1$ (monotone convergence).
\end{proof}
No assumption about the boundary of $M$ or about non-hyperbolic interior components enters.

\paragraph{Coefficients.} With Taylor coefficients $\lambda(c^*+e)=\lambda_1e+\lambda_2e^2+\lambda_3e^3+O(e^4)$, series reversion gives $a_1=\lambda_1^{-1}$, $a_2=-\lambda_2\lambda_1^{-3}$, $a_3=(2\lambda_2^2-\lambda_1\lambda_3)\lambda_1^{-5}$. Writing $g_p(c)=f_c^p(0)$ and differentiating the fixed-point equation at the centre (where the return map has zero derivative in $z$) gives $\lambda_1=2g_p'(c^*)\prod_{j=1}^{p-1}2f^j_{c^*}(0)$. Higher coefficients are obtained from finite power-series computations that are exact as formal series: the generator iterates the return map on jets (each pass gains one order, since $F_z(0,0)=0$); one verifier takes two Newton steps on truncated series (orders $1\to2\to4$); the other solves the system $(f_c^p(z)-z,\ \prod_j 2f_c^j(z)-\lambda)$ order by order in $\lambda$. Evaluated over a certified disk containing $c^*$, these enclose the values at the exact centre. For the cardioid and the period-2 disk, $\psi_1(\lambda)=\lambda/2-\lambda^2/4$ and $\psi_2(\lambda)=-1+\lambda/4$, giving the exact areas $3\pi/8$ and $\pi/16$.

\paragraph{Certification of each component.} For each candidate centre (found by floating-point Newton iteration):
\begin{enumerate}
\item a contraction test on a closed disk $D$ with a fixed preconditioner $Y$ ($\sup_D|1-Yg_p'|\le q<1$ and $|Yg_p(t)|+qR\le R$) proves a unique zero $c^*$ of $g_p$ in $D$;
\item for every proper divisor $m$ of $p$, $0\notin f^m(0)$ is checked in ball arithmetic over the box, so the critical point has exact period $p$ and $c^*$ is the centre of a hyperbolic component of period $p$;
\item lower endpoints of $|a_n|$ are clamped at zero before squaring, and the area terms are summed with outward rounding.
\end{enumerate}
Components of different periods are distinct. Within a period, the stored boxes are pairwise disjoint (checked globally by exact rational comparison of all squares), so centres and components are distinct. Boxes strictly above the real axis count twice (the conjugate component); boxes on the axis count once and contain a real centre. Missing components only weaken the bound.

\paragraph{Data.} The certificate has 1,679,806 rows in twelve files (Table~\ref{tab:lb}): a base set covering periods 3--504 (443 periods present; 508 was the search cap) and eleven supplements restricted to periods 13--136 that exclude the base set's boxes. The file names of the supplements begin with ``prim''; this is a label, not a claim that the components are primitive. The audited sums use three coefficients ($a_1,a_2,a_3$) for all supplement rows and for base rows of period at most 64, and two coefficients for base rows above period 64.

\begin{table}[h]
\centering
\begin{tabular}{@{}lrrr@{}}
\toprule
Data set & Periods & Rows & Components \\
\midrule
Periods 1--2 (exact, $7\pi/16$) & 1--2 & -- & 2 \\
Base & 3--504 & 187,585 & 374,785 \\
Supplements & 13--136 & 1,492,221 & 2,977,807 \\
\midrule
Total & & 1,679,806 & 3,352,594 \\
\bottomrule
\end{tabular}
\caption{Lower-bound certificate. ``Components'' counts conjugates. Cumulative bound after the base set: $1.5063005$; after all supplements: $1.50648354$ (independent audit's sum), stated conservatively as $1.5064820$.}
\label{tab:lb}
\end{table}

\section{The upper bound: a certified quadtree}

$M$ lies in the closed disk of radius 2 and is symmetric under conjugation. We cover the upper half by the squares $[-2,0]\times[0,2]$ and $[0,2]\times[0,2]$ (area 8) and subdivide dyadically. If $A_E$ is the total area credited to cells certified disjoint from $M$ (the credited cells are pairwise disjoint and inside the root; unresolved cells are not stored, and no tiling is needed), then $\mu(M)\le2(8-A_E)$; the real axis has measure zero. Write $z_1=c$, $z_{k+1}=z_k^2+c$, $z_k'=dz_k/dc$, and let $\Phi$ be the B\"ottcher map of $\mathbb{C}\setminus M$, $\Phi(c)\sim c$ at infinity, with $G=\log|\Phi|=\lim_k2^{1-k}\log|z_k|$. Below, $|\nabla G|=|\Phi'/\Phi|$ (not the Wirtinger derivative $|\partial G/\partial c|$, which is half of it).

\begin{lemma}[escape]
If $|z_n|>2$ for some $n$ then $c\notin M$. A cell whose ball orbit satisfies this at some step is exterior.
\end{lemma}

\begin{lemma}[potential and gradient]\label{lem:tail}
Let $N\ge1$, $|z_N|\ge4$, $\delta=|z_N|^{-2}$, $\varepsilon=|c|\delta\le1/8$, $\nu=2^{1-N}$. Then $c\notin M$ and
\[
  \big|G(c)-\nu\log|z_N|\big|\le\nu\varepsilon,\qquad
  |\nabla G(c)|\le\nu\Big(\frac{|z_N'|}{|z_N|}(1+3\varepsilon)+\delta\Big).
\]
\end{lemma}
\begin{proof}
With $\delta_k=|z_k|^{-2}$ and $\varepsilon_k=|c|\delta_k$ for $k\ge N$: $|z_{k+1}|\ge|z_k|^2(1-\varepsilon_k)\ge\frac78|z_k|^2$, hence $\delta_{k+1}\le\frac{64}{49}\delta_k^2\le\delta_k/12$ and $\varepsilon_{k+1}\le\varepsilon_k/12$. Since $|\log|1+u||\le|u|/(1-|u|)$, telescoping $G-\nu\log|z_N|=\sum_{k\ge N}2^{-k}\log|1+c/z_k^2|$ gives $|G-\nu\log|z_N||\le2^{-N}\frac87\varepsilon\sum_j24^{-j}=\nu\frac{96}{161}\varepsilon$. For the gradient, $D_k=2^{1-k}z_k'/z_k$ satisfies $D_{k+1}=(D_k+2^{-k}z_k^{-2})/(1+c/z_k^2)$. With $P=\prod_{k\ge N}(1-\varepsilon_k)^{-1}\le(1-\frac{12}{11}\varepsilon)^{-1}\le\frac{22}{19}$, $P-1\le\frac{24}{19}\varepsilon$ and $\sum_{k\ge N}2^{-k}\delta_k\le\nu\frac{12}{23}\delta$, one gets $|D_\infty|\le(1+\frac{24}{19}\varepsilon)|D_N|+\nu\frac{264}{437}\delta$, which implies the stated bound. On a small simply connected neighbourhood of $c$ fix compatible branches of the logarithms; the tail series converges locally uniformly there by the geometric decay, so the limit is analytic, equals $\log\Phi$ up to a constant, and may be differentiated termwise, giving $D_\infty=\Phi'/\Phi$.
\end{proof}

\begin{lemma}[distance]\label{lem:koebe}
For $c\notin M$, $\operatorname{dist}(c,M)\ge\big(1-e^{-2G(c)}\big)/\big(4|\nabla G(c)|\big)$.
\end{lemma}
\begin{proof}
Let $\Psi=\Phi^{-1}$, $w_0=\Phi(c)$, $r=|w_0|>1$, and $S$ a disk automorphism with $S(0)=1/w_0$. Then $F=\Psi\circ(1/S)$ is injective and holomorphic on the disk of radius $1/r$ (its pole lies on the boundary), with $|F'(0)|=(r^2-1)/|\Phi'(c)|$. Koebe's quarter theorem \cite{duren1983} on that disk gives $\operatorname{dist}(c,M)\ge(r^2-1)/(4r|\Phi'(c)|)$, which equals the stated expression since $r=e^{G}$.
\end{proof}

A cell of side $s$ with centre $c_0$ is exterior when the bound of Lemma~\ref{lem:koebe} at $c_0$, computed from a lower bound on $G$ and an upper bound on $|\nabla G|$ (Lemma~\ref{lem:tail}, applied only at $c_0$ itself), exceeds the cell's circumradius $s/\sqrt2$. The final checker requires the stronger hypotheses $|z_N|\ge64$ and $\varepsilon\le10^{-3}$ and uses $1-e^{-2G}=-\operatorname{expm1}(-2G)$ to avoid cancellation. At the deepest level a cell may be credited in part: if the certified distance at the centre is $d_0$, the cell's exterior contains a disk of radius $\min(d_0,s/2)$ and an axis-parallel square inscribed in the disk of radius $\min(d_0,s/\sqrt2)$, so a credit of $m\,s^2/256$ is accepted when it does not exceed the larger of $\pi\min(d_0,s/2)^2$ and $2\min(d_0,s/\sqrt2)^2$.

\paragraph{Data.} 808,356,698 credited cells (442,614,788 full exterior, 365,741,910 partial) in 4,242 files, base depth 17 with two further exterior levels (maximum depth 19), 14,537 occupied depth-9 roots. The exact credit is $C=126977322011966$ in units of $(2\cdot2^{-19})^2/256$, so $\mu(M)\le16-8C/(256\cdot4^{19})=6880083171681/2^{42}$.

\section{Verification}

Each certificate was produced by one program and audited by others written from the stated lemmas and the certificate format, without reading the generator's code.
\begin{itemize}
\item \textbf{Lower bound.} Two separately written auditors re-established every root (contraction on a disk), exact period, the coefficients and global disjointness of all 1,679,806 boxes, at adaptive precision of 192--1536 bits and 512--2048 bits respectively, and both reject corrupted certificates. The second summed the certified terms to $1.50648354\ldots$; the first, organised as a chain of checks over the twelve files, reaches $1.5064821\ldots$.
\item \textbf{Upper bound.} A separately written auditor re-verified every one of the 808,356,698 credited cells in Arb at adaptive precision (192, 384, 768 and 1536 bits), checked disjointness and containment by exact dyadic (Morton) intervals, and recomputed the exact credit by integer arithmetic (about 4.1 hours on four cores). It rejects corrupted certificates. A first attempt at fixed 192 bits stopped on an enclosure failure caused by cancellation in $1-e^{-2G}$; the failure was a conservative rejection, and the cell passes at higher precision.
\item \textbf{Review.} An independent reviewer checked the lemmas above (the tail estimate and the coefficient formulas are given here as the reviewer supplied or confirmed them), recomputed all upper-bound geometry and credits from the raw records with separate integer code, checked global separation of all lower-bound squares exactly, reconstructed the lower-bound sum from the audit's certified terms, and checked the coefficient constructions in exact arithmetic for periods 1--3. The reviewer did not re-run the Arb witness checks.
\end{itemize}

\paragraph{What is trusted, and what is shared.} The auditors share the input candidates (as untrusted data), the theorems used (uniformization of hyperbolic components and the area identity; escape, B\"ottcher uniformization, the tail estimates and Koebe's theorem) and the arithmetic libraries (Arb/FLINT via python-flint, NumPy for parsing, Python). Separately written code does not make the mathematics independent. Model-vendor diversity among the agents that wrote the code is not independence of the theorems either.

\paragraph{Reproducing the checks.} The auditors save resumable checkpoints. A checkpoint records the result of a completed check and is not re-verified when reused, so a checkpoint must never be treated as part of the proof: verification must start in a fresh, empty working directory. The certificates, the auditors and a clean-run script with pinned software versions are deposited with this version (see Data).

\section{How this result was produced}

This work was done with an automated research harness in which teams of AI agents (models from Anthropic and OpenAI) work in rounds: a theorist states lemmas, researchers write generators and checkers, a skeptic attacks claims, and an editor tracks what is established. The author chose the problem, set the direction and the acceptance rule (a bound counts only when a separately written checker passes the whole certificate), approved long computations, and checked the comparison with prior work. The lemmas, code and certificates were produced by the agents over two runs on 9 and 10 November 2026 on a laptop; the review was done by a further agent that had not seen the runs. The author has not checked every detail of the proofs and code himself; the published certificates and checkers are meant to make that possible for others.

\section*{Data}

Certificates, auditors and the clean-run script: see the data record linked from this paper's page.

\section*{Revision note}

Version 1 (10 November 2026) stated that the lower bound exceeded every published lower bound we knew of; a stronger ball-arithmetic lower bound \cite{sciencejournal2026} had appeared in October 2026, and that statement is withdrawn. Version 2 also corrects the counts (3,352,594 components including periods 1--2; base periods 3--504), the precision ranges of the two lower-bound audits, and the description of which coefficients are summed; drops the word ``primitive''; adds proofs of the area identity and the tail estimate; adds the warning on checkpoints; and removes the claim to be the first certified bounds of this quality.

\bibliographystyle{plain}
\bibliography{refs}

\end{document}
